\documentclass[a4paper]{article}     
\usepackage{amsfonts}
\usepackage{amsmath}
\usepackage{amssymb}
\usepackage{graphicx}
\usepackage{cancel}

\newcommand{\asa}{\Leftrightarrow} %als slechts als
\newcommand{\adR}{\Rightarrow} %als dan Rechts
\newcommand{\adL}{\Leftarrow}
\newcommand{\Lh}{\left(} %Links haakje
\newcommand{\Rh}{\right)}
\newcommand{\Lv}{\left[} %Links vierkanthaakje
\newcommand{\Rv}{\right]}
\newcommand{\La}{\left\{} %Links acolade
\newcommand{\Ra}{\right\}}
\newcommand{\Ras}{\right.} %Rechts acolade stelsel (omdat om stelsels te maken heb je een punt nodig
\newcommand{\pnR}{\rightarrow} %pijl naar Rechts
\newcommand{\limiet}{\lim\limits_}
\newcommand{\schrap}{\cancel}
\newcommand{\schrapb}{\bcancel} %backslash
\newcommand{\schrapk}{\xcancel} %kruis
\newcommand{\vervang}{\cancelto}
\newcommand{\intgrl}{\displaystyle\int} %integraal teken (bij het berekenen van primitieven)

\begin{document}

\noindent \large Auteur: Yoren Collier \\
\noindent \large Studierichting: Informatica\\
\noindent \large Oefeningenreeks 4A nr. 38.2\\

\medskip

\normalsize

$ I =  \intgrl \dfrac{2x+11}{x^2 +6x +10} dx = 2\intgrl \dfrac{x}{x^2 +6x +10} dx + 11\intgrl \dfrac{1}{\Lh x+3 \Rh^2 +1} dx$\\

$= 2 \intgrl \Lh \dfrac{2x +6}{2 \Lh x^2 +6x +10 \Rh} - \dfrac{3}{x^2 +6x+ 10} \Rh dx + 11 \intgrl \dfrac{1}{\Lh x+3 \Rh^2 +1} dx$\\

Stel: $t = x^2 +6x +10 \adR dt = 2x +6 dx$ en stel: $y = x +3 \adR dy = dx$\\

$= 2 \Lv \intgrl \dfrac{1}{2 \Lh x^2 +6x +10 \Rh} \cdot \Lh 2x +6 dx \Rh - \intgrl \dfrac{3}{\Lh x+3 \Rh^2 +1 } dx\Rv +11 \intgrl \dfrac{1}{\Lh x+3 \Rh^2 +1 } dx$\\

$= 2 \Lv \intgrl \dfrac{1}{2t} dt - \intgrl \dfrac{3}{y^2 +1} dy\Rv +11 \intgrl \dfrac{1}{y^2 +1} dy$\\

$= 2 \Lv \dfrac{1}{2} \ln |t| +c -3\operatorname{Bgtan} \Lh y \Rh +c \Rv +11 \operatorname{Bgtan} \Lh y \Rh +c$\\

$= \ln |x^2 +6x +10| -6 \operatorname{Bgtan} \Lh x+3 \Rh +11 \operatorname{Bgtan} \Lh x+3 \Rh +c$\\

$= \ln |x^2 +6x +10| +5 \operatorname{Bgtan} \Lh x+3 \Rh +c$\\

\begin{thebibliography}{999}
%\bibitem{a} Referentie
\end{thebibliography}
\end{document} 